\section{Conclusiones}

La práctica permitió conocer el funcionamiento de GitHub Issues tanto desde la interfaz web como desde GitHub CLI. Se comprobó que la herramienta facilita la organización de actividades dentro de un proyecto colaborativo mediante la creación, visualización, asignación y seguimiento de incidencias.

Asimismo, el uso de GitHub CLI permitió realizar las mismas tareas directamente desde la terminal, agilizando la administración de Issues y favoreciendo el trabajo colaborativo entre los integrantes del equipo. Finalmente, se logró simular un flujo de trabajo real donde las incidencias pueden ser registradas, asignadas y monitoreadas hasta su resolución.

Además, se realizó el seguimiento y cierre de los Issues asignados, documentando el proceso mediante comentarios, ramas de trabajo, commits y actualizaciones en el repositorio, lo que permitió completar el ciclo de gestión de incidencias dentro del proyecto.